\documentclass[10pt,a4paper]{amsart}
\usepackage{amsmath,amssymb,amsthm,graphicx}
% Retained minimal preamble for the Stacks Project chapter more-algebra.tex (More on Algebra).
% The chapter file pinned-original-source.tex carries only \input{preamble} on its first line, so no
% preamble text lives inside the chapter itself.  The Stacks Project uses one single shared preamble for
% every chapter; the snapshot stages it identically as sources/stacks-chow/preamble.tex and
% sources/stacks-etale-cohomology/preamble.tex.  The block below is copied mechanically, line for line,
% from sources/stacks-chow/preamble.tex lines 14-18 (the Xy-pic setup) and 154-217 (the native theorem
% declarations together with the mathematical macros).  Nothing was added, removed or redefined; the
% wrapper class and the amsmath/amssymb/amsthm/graphicx base packages are supplied by the assembler.
% Chapter: more-algebra.tex -- More on Algebra.
% For commutative diagrams we use Xy-pic
\usepackage[all]{xy}

% We use 2cell for 2-commutative diagrams.
\xyoption{2cell}

% Theorem environments.
%
\theoremstyle{plain}
\newtheorem{theorem}[subsection]{Theorem}
\newtheorem{proposition}[subsection]{Proposition}
\newtheorem{lemma}[subsection]{Lemma}

\theoremstyle{definition}
\newtheorem{definition}[subsection]{Definition}
\newtheorem{example}[subsection]{Example}
\newtheorem{exercise}[subsection]{Exercise}
\newtheorem{situation}[subsection]{Situation}

\theoremstyle{remark}
\newtheorem{remark}[subsection]{Remark}
\newtheorem{remarks}[subsection]{Remarks}

\numberwithin{equation}{subsection}

% Macros
%
\def\lim{\mathop{\mathrm{lim}}\nolimits}
\def\colim{\mathop{\mathrm{colim}}\nolimits}
\def\Spec{\mathop{\mathrm{Spec}}}
\def\Hom{\mathop{\mathrm{Hom}}\nolimits}
\def\Ext{\mathop{\mathrm{Ext}}\nolimits}
\def\SheafHom{\mathop{\mathcal{H}\!\mathit{om}}\nolimits}
\def\SheafExt{\mathop{\mathcal{E}\!\mathit{xt}}\nolimits}
\def\Sch{\mathit{Sch}}
\def\Mor{\mathop{\mathrm{Mor}}\nolimits}
\def\Ob{\mathop{\mathrm{Ob}}\nolimits}
\def\Sh{\mathop{\mathit{Sh}}\nolimits}
\def\NL{\mathop{N\!L}\nolimits}
\def\CH{\mathop{\mathrm{CH}}\nolimits}
\def\proetale{{pro\text{-}\acute{e}tale}}
\def\etale{{\acute{e}tale}}
\def\QCoh{\mathit{QCoh}}
\def\Ker{\mathop{\mathrm{Ker}}}
\def\Im{\mathop{\mathrm{Im}}}
\def\Coker{\mathop{\mathrm{Coker}}}
\def\Coim{\mathop{\mathrm{Coim}}}

% Boxtimes
%
\DeclareMathSymbol{\boxtimes}{\mathbin}{AMSa}{"02}

%
% Macros for moduli stacks/spaces
%
\def\QCohstack{\mathcal{QC}\!\mathit{oh}}
\def\Cohstack{\mathcal{C}\!\mathit{oh}}
\def\Spacesstack{\mathcal{S}\!\mathit{paces}}
\def\Quotfunctor{\mathrm{Quot}}
\def\Hilbfunctor{\mathrm{Hilb}}
\def\Curvesstack{\mathcal{C}\!\mathit{urves}}
\def\Polarizedstack{\mathcal{P}\!\mathit{olarized}}
\def\Complexesstack{\mathcal{C}\!\mathit{omplexes}}
% \Pic is the operator that assigns to X its picard group, usage \Pic(X)
% \Picardstack_{X/B} denotes the Picard stack of X over B
% \Picardfunctor_{X/B} denotes the Picard functor of X over B
\def\Pic{\mathop{\mathrm{Pic}}\nolimits}
\def\Picardstack{\mathcal{P}\!\mathit{ic}}
\def\Picardfunctor{\mathrm{Pic}}
\def\Deformationcategory{\mathcal{D}\!\mathit{ef}}
\usepackage{cleveref}
\crefname{originallemma}{original lemma}{original lemmas}
\Crefname{originallemma}{Original lemma}{Original lemmas}
\makeatletter
\newlabel{more-algebra-lemma-compose}{{0A5Y}{source}{}{sourcepointer-1}{}}
\newlabel{more-algebra-lemma-compose@cref}{{[originallemma][0][]0A5Y}{[1][0][]source}{}{}{}}
\makeatother
\crefname{originallemma}{original lemma}{original lemmas}
\Crefname{originallemma}{Original lemma}{Original lemmas}
\makeatletter
\newlabel{homology-definition-associated-simple-complex}{{012Z}{source}{}{sourcepointer-2}{}}
\newlabel{homology-definition-associated-simple-complex@cref}{{[originallemma][0][]012Z}{[1][0][]source}{}{}{}}
\makeatother
\title{The internal Hom adjunction for complexes as an isomorphism of Hom complexes}
\author{The Stacks Project Authors}
\date{stacks-project:more-algebra.tex:local-snapshot-20261012; GNU FDL 1.2 or later}
\begin{document}
\maketitle
\setcounter{subsection}{1}
\noindent\textit{Editorial note.} All delivered mathematical text is the author's own, copied from the pinned Stacks Project chapter file more-algebra.tex; no statement or proof step was rewritten, reordered or removed.\par\medskip
\noindent\textit{Editorial note.} The chapter file pinned-original-source.tex carries only \textbackslash{}input\{preamble\} on its first line, so no preamble text lives inside the chapter. The retained file sources/stacks-more-algebra/retained-preamble.tex is a mechanical, line-for-line copy of the project's single shared preamble as staged in the snapshot (Xy-pic setup plus theorem declarations and mathematical macros). Nothing was added, removed or redefined, and the excerpt serves only as preamble, never as a span coordinate.\par\medskip
\noindent\textit{Editorial note.} Both counted spans start at the \textbackslash{}begin line of their environment, so the statement and the author's proof are delivered whole; surrounding prose of the section is deliberately outside the counted spans and is carried as standard prerequisites instead.\par\medskip
\noindent\textit{Editorial note.} The citation fields (source\_id, source\_version, authors, source\_url, license.name, source\_title) are identical to those of every other card this lane has produced from this chapter, so the batch selector treats them as one source.\par\medskip
\noindent\textit{Allowed standard prerequisites (editorial).}\begin{itemize}
\item Internal Hom complex of two complexes and the total tensor product of complexes. constructed and discussed in the same section, source lines 17800-17879, before the counted statement. Supplies both sides of the isomorphism and their differentials.
\item The tensor-Hom identification in each bidegree, viewing a map into an internal Hom as a map out of a tensor product. used at proof lines 17895-17905 under the displayed equation. Produces the components of the right hand side from those of the left.
\item The sign rule for the differential of an associated simple complex. the author cites Homology, Definition homology-definition-associated-simple-complex at proof lines 17930-17933. Fixes the signs with which the two differentials are compared.
\end{itemize}\medskip
\ifdefined\hypertarget\hypertarget{sourcepointer-1}{}\fi
\noindent\textit{Source reference.} Original lemma 0A5Y: focused-seeds/AGP/sources/stacks-snapshot-20261011/tags.txt, source lines 12826--12826. The lemma's own canonical tag, quoted so that readers can resolve the label inside the delivered text; it points at the counted statement itself and carries no proof content.\par\medskip
\ifdefined\hypertarget\hypertarget{sourcepointer-2}{}\fi
\noindent\textit{Source reference.} Original lemma 012Z: focused-seeds/AGP/sources/stacks-snapshot-20261011/tags.txt, source lines 1417--1417. Navigation pointer to a statement the counted proof refers to but does not contain; nothing load-bearing is delegated to it.\par\medskip
\begin{lemma}
\label{lemma-compose}
Let $R$ be a ring. Given complexes $K^\bullet, L^\bullet, M^\bullet$
of $R$-modules there is a canonical isomorphism
$$
\Hom^\bullet(K^\bullet, \Hom^\bullet(L^\bullet, M^\bullet))
=
\Hom^\bullet(\text{Tot}(K^\bullet \otimes_R L^\bullet), M^\bullet)
$$
of complexes of $R$-modules.
\end{lemma}

\begin{proof}
Let $\alpha$ be an element of degree $n$ on the left hand side.
Thus
$$
\alpha = (\alpha^{p, q}) \in
\prod\nolimits_{p + q = n} \Hom_R(K^{-q}, \Hom^p(L^\bullet, M^\bullet))
$$
Each $\alpha^{p, q}$ is an element
$$
\alpha^{p, q} = (\alpha^{r, s, q}) \in
\prod\nolimits_{r + s + q = n} \Hom_R(K^{-q}, \Hom_R(L^{-s}, M^r))
$$
If we make the identifications
\begin{equation}
\label{equation-identification}
\Hom_R(K^{-q}, \Hom_R(L^{-s}, M^r)) = \Hom_R(K^{-q} \otimes_R L^{-s}, M^r)
\end{equation}
then by our sign rules we get
\begin{align*}
\text{d}(\alpha^{r, s, q})
& =
\text{d}_{\Hom^\bullet(L^\bullet, M^\bullet)} \circ \alpha^{r, s, q}
- (-1)^n \alpha^{r, s, q} \circ \text{d}_K \\
& =
\text{d}_M \circ \alpha^{r, s, q}
- (-1)^{r + s} \alpha^{r, s, q} \circ \text{d}_L
- (-1)^{r + s + q} \alpha^{r, s, q} \circ \text{d}_K
\end{align*}
On the other hand, if $\beta$ is an element of degree $n$ of
the right hand side, then
$$
\beta = (\beta^{r, s, q}) \in \prod\nolimits_{r + s + q = n}
\Hom_R(K^{-q} \otimes_R L^{-s}, M^r)
$$
and by our sign rule (Homology, Definition
\ref{homology-definition-associated-simple-complex}) we get
\begin{align*}
\text{d}(\beta^{r, s, q})
& =
\text{d}_M \circ \beta^{r, s, q}
- (-1)^n \beta^{r, s, q} \circ
\text{d}_{\text{Tot}(K^\bullet \otimes L^\bullet)} \\
& =
\text{d}_M \circ \beta^{r, s, q}
- (-1)^{r + s + q} \left(
\beta^{r, s, q} \circ \text{d}_K + (-1)^{-q} \beta^{r, s, q} \circ \text{d}_L
\right)
\end{align*}
Thus we see that the map induced by the identifications
(\ref{equation-identification}) indeed is a morphism of complexes.
\end{proof}

\end{document}
